\subsection{\texorpdfstring{空间 \(E_{A}\)（A 有界）}{空间 E\_\{A\}（A 有界）}}

设 E 是一个局部凸空间，A 是 E 的一个凸均衡子集。我们回忆，\(E_{A}\) 表示由 A 生成的向量空间，并以 A 的度规 \(p_{A}\) 作为半范数（II, 第 26 页, 例 3）。我们立即验证，\(E_{A}\) 到 E 中的典范单射连续，当且仅当 A 有界。此外，若 E 是 Hausdorff 的，则有界集 A 不包含任何直线（III, 第 2 页, 注 4），于是 \(p_{A}\) 是一个范数（II, loc. cit.）。

我们称一致空间 X 是半完备的，如果 X 中每个 Cauchy 序列都收敛。完备的一致空间是半完备的；但其逆并不总成立（GT, II, § 4, 习题 4）；不过，可度量化的半完备空间是完备的（GT, IX, § 2, No.~6, 命题 9）。

命题 8. --- 设 E 是一个 Hausdorff 局部凸空间，A 是 E 的一个闭的、均衡的、有界的且凸的子集。设 \((x_{n})\) 是 \(E_{A}\) 中的一个 Cauchy 序列。则此序列在 \(E_{A}\) 中收敛，当且仅当它在 E 中收敛。

\(E_{A}\) 到 E 中的典范单射是连续的。因此，若 \((x_{n})\) 在 \(E_{A}\) 中收敛，则它在 E 中收敛。反之，设 \((x_{n})\) 在 E 中收敛于 y。存在一个递增的整数序列 \((n_{k})\)，使得 \(p_{\mathbf{A}}(x_{m}-x_{n}) \leqslant 2^{-k-1}\) 成立，只要 \(m \geqslant n_{k}\) 且 \(n \geqslant n_{k}\)。于是序列 \((x_{n_{k}} + 2^{-k}\mathbf{A})\) 是递减的。由于 A 在 E 中是闭的，我们有 \(y \in \bigcap_{k} (x_{n_{k}} + 2^{-k}\mathbf{A})\)，这就证明了 \((x_{n_{k}})\)，从而 \((x_{n})\)，在 \(E_{A}\) 中收敛于 y。

推论. --- 若 A 是半完备的（特别是，完备的），则 \(E_{A}\) 是一个 Banach 空间。

事实上，\(E_{A}\) 中的一个 Cauchy 序列对 E 的拓扑而言也是 Cauchy 序列，并且含于一个与 A 位似的集合中，因此在 E 中收敛。

